\section{A GitHub account and access to the repository}
\label{sec:github}

\subsection{Background}

\term{Git} is a program that records the history of a folder of files: every saved state (a
\term{commit}) can be recalled, compared and shared. A folder managed by Git is called a
\term{repository}. \term{GitHub} is a website that stores repositories on the internet so that
several people can work on the same one. The team's repository is

\begin{center}
\repourl
\end{center}

and it belongs to the project lead, Andrew Shahvosian, whose GitHub user name is
\texttt{ashahvos}. To read the repository you may need an account that the lead has invited; to
contribute to it you certainly do. The whole section is done in a web browser (Edge, Chrome or
Firefox); nothing is installed yet.

\subsection{Create the account}

Skip to Step~\ref{st:sendname} if you already have a GitHub account and can sign in to it.

\step{Open the sign-up page}
Open the browser, click in the address bar at the top, type the following address and press
\key{Enter}.
\begin{center}\url{https://github.com/signup}\end{center}
\checkpoint{A page headed \emph{Sign up for GitHub} (or similar) asks for an email address.}

\step{Fill in the form}
The layout of the page changes from time to time, but it always asks for the same four items.
\begin{enumerate}
\item \textbf{Email.} Enter an address whose inbox you can open right now; your USC address is a
  good choice.
\item \textbf{Password.} GitHub accepts a password of at least 15 characters, or of at least 8
  characters that include a number and a lowercase letter. Do not reuse a password from another
  site, and store it in a password manager or write it down now.
\item \textbf{Username.} This name is public and appears on every contribution you make. Use
  letters, digits and single hyphens only. A recognizable form of your own name is best, for
  example the part of your USC address before the \code|@| sign.
\item \textbf{Country or region.} Select it from the list.
\end{enumerate}
Leave any box about product updates unticked unless you want the emails, then click the button
at the bottom of the form (\ui{Create account} or \ui{Continue}).

\step{Prove that you are a person and confirm the email address}
GitHub shows a short visual puzzle; solve it. It then sends a numeric code to the email address
you gave. Open your inbox in another browser tab, find the message from GitHub (check the spam
folder if it has not arrived after two minutes), and type the code into the sign-up page.
\checkpoint{The browser shows your GitHub home page, or a sign-in page that accepts your new
user name and password. Any questions about your interests can be skipped with the link
\ui{Skip personalization} (or \ui{Skip this for now}); the free plan is sufficient.}

\begin{note}[Two-factor authentication]
GitHub strongly recommends, and for people who contribute code may require, a second sign-in
factor such as an authenticator app on your phone. If GitHub asks you to set it up, follow its
instructions; the request is genuine, and it does not interfere with anything in this guide.
\end{note}

\subsection{Obtain access}

\step{Send your user name to the project lead}\label{st:sendname}
Send Andrew a message through the team's usual channel that contains your GitHub user name
exactly as you typed it, including any hyphens. Then wait for his reply; he carries out
Step~\ref{st:invite}.
\checkpoint{The lead has confirmed that the invitation was sent.}

\step{For the project lead only: invite the team member}\label{st:invite}
\begin{enumerate}
\item Sign in to GitHub and open the page of the repository,\\ \repourl.
\item Click the \ui{Settings} tab of the repository (the rightmost tab under the repository
  name; if it is hidden, open the \ui{\ldots} menu at the end of the row of tabs).
\item In the left-hand list, under \ui{Access}, click \ui{Collaborators}. GitHub may ask for
  your password again.
\item Click \ui{Add people}, type the team member's user name, select the matching account in
  the list, and click the green button that names the account and the repository.
\end{enumerate}
The same page shows whether the repository is \emph{Private} or \emph{Public}. A private
repository cannot be seen at all by an account that has not accepted an invitation.
\checkpoint{The team member's name is listed on the Collaborators page with the remark
\emph{Pending invite}.}

\step{Accept the invitation}
GitHub sends you an email which says that \texttt{ashahvos} has invited you to collaborate on
the repository \texttt{ashahvos/camera-to-latex}. Open it and click \ui{View invitation}, then
\ui{Accept invitation}.
If the email cannot be found, sign in to GitHub and open this address instead:
\begin{center}\url{https://github.com/ashahvos/camera-to-latex/invitations}\end{center}
An invitation expires after seven days; if yours has expired, ask the lead to send it again.
\checkpoint{The browser shows the repository page: a list of folders and files
(\file{firmware}, \file{model}, \file{pipeline}, \file{tools}, \ldots) and, below it, the text of
the file \file{README.md} under the heading \emph{Camera to LaTeX}.}

\subsection{A first look at the repository}

\step{Read a file in the browser}
Nothing needs to be installed to read the project. On the repository page, click the folder
\file{docs}, then a PDF file in it; GitHub displays the document, and the download button
(an arrow pointing into a tray, at the top right of the file view) saves a copy. Click the
repository name \ui{camera-to-latex} at the top of the page to return to the start. The same
method works for \file{sheets/sheets.pdf} and \file{test\_set/test\_lines.pdf}, the two documents
that will later be printed.
\checkpoint{You have opened one file of the repository and returned to its first page.}

Table~\ref{tab:repo} lists what the repository contains. The names will become familiar in the
quick start of Section~\ref{sec:quick}.

\begin{table}[H]
\centering
\begin{tabular}{L{3.5cm}L{8.1cm}}
\toprule
\multicolumn{1}{c}{\hd{Name}} & \multicolumn{1}{c}{\hd{Contents}} \\
\midrule
\file{README.md} & A one-page summary of the project and its commands \\
\file{config.py} & Every adjustable number of the pipeline, in one file \\
\file{symbols.csv} & The 159 classes (182 symbols) the classifier knows \\
\file{requirements.txt} & The Python packages the project needs \\
\file{firmware/capture/} & The program (``sketch'') that runs on the board \\
\file{pipeline/} & Capture, segmentation, classification, layout, PDF output and the application \\
\file{model/} & Generation of training images, training and evaluation of the classifier \\
\file{tools/} & Printable sheets, test lines, evaluation and review of results \\
\file{sheets/} & The printable sheets of symbols (\file{sheets.pdf}, 9 pages) \\
\file{test\_set/} & The 80 test lines (\file{lines.tsv}, \file{test\_lines.pdf}) and their photographs \\
\file{tests/} & Automatic tests of the code \\
\file{docs/} & This guide and its companion \\
\bottomrule
\end{tabular}
\caption{The contents of the repository.}
\label{tab:repo}
\end{table}
